\documentclass[11pt,a4paper]{article}

\usepackage[utf8]{inputenc}
\usepackage[T1]{fontenc}
\usepackage{lmodern}
\usepackage[margin=2.5cm]{geometry}
\usepackage{amsmath,amssymb}
\usepackage{booktabs,tabularx,array}
\usepackage{xcolor}
\usepackage{tikz}
\usetikzlibrary{positioning,arrows.meta,fit,backgrounds,shapes.geometric,calc}
\usepackage{listings}
\usepackage{enumitem}
\usepackage{caption}
\usepackage{float}
\usepackage{fancyhdr}
\usepackage[hidelinks]{hyperref}

% --- Nombres en español (sin babel) ---------------------------------------
\renewcommand{\contentsname}{Índice}
\renewcommand{\tablename}{Tabla}
\renewcommand{\figurename}{Figura}
\renewcommand{\refname}{Referencias}
\renewcommand{\lstlistingname}{Código}

\definecolor{azul}{HTML}{1F4E79}
\definecolor{verde}{HTML}{2E7D32}
\definecolor{rojo}{HTML}{C62828}
\definecolor{naranja}{HTML}{EF6C00}
\definecolor{gris}{HTML}{F4F6F8}

\lstset{
  language=Python, basicstyle=\ttfamily\footnotesize, keywordstyle=\color{azul}\bfseries,
  commentstyle=\color{gray}\itshape, stringstyle=\color{verde}, backgroundcolor=\color{gris},
  frame=single, rulecolor=\color{gray!40}, breaklines=true, numbers=left,
  numberstyle=\tiny\color{gray}, showstringspaces=false, inputencoding=utf8, extendedchars=true,
  literate={á}{{\'a}}1 {é}{{\'e}}1 {í}{{\'i}}1 {ó}{{\'o}}1 {ú}{{\'u}}1 {ñ}{{\~n}}1
           {Á}{{\'A}}1 {É}{{\'E}}1 {Í}{{\'I}}1 {Ó}{{\'O}}1 {Ú}{{\'U}}1
}

\pagestyle{fancy}
\fancyhf{}
\lhead{\small Agente DevOps de revisión de Pull Requests}
\rhead{\small Tecnology Trend}
\cfoot{\thepage}

\newcommand{\Pdef}{P_{\text{def}}}

\begin{document}
\sloppy

% ============================== PORTADA ===================================
\begin{titlepage}
\centering
{\large FUNDACIÓN UNIVERSITARIA LOS LIBERTADORES\par}
{\large Facultad de Ingeniería y Ciencias Básicas\par}
{\large Ingeniería de Sistemas\par}
\vspace{3cm}
{\Huge\bfseries Agente Inteligente de Revisión de Código y Gestión de Pull Requests\par}
\vspace{0.6cm}
{\Large Proyecto 4 --- DevOps\par}
\vspace{0.4cm}
{\large Taller de Ingeniería de Agentes Inteligentes\par}
\vspace{0.3cm}
{\large\itshape Entrega 3: Implementación, pruebas y conclusiones\par}
\vspace{3cm}
{\large\textbf{Presentado por:}\\ Simon Hurtado\par}
\vspace{0.8cm}
{\large\textbf{Docente:}\\ Profesor Gonzalo\par}
\vspace{0.8cm}
{\large\textbf{Asignatura:}\\ Tecnology Trend\par}
\vfill
{\large Bogotá D.C., noviembre de 2026\par}
\end{titlepage}

\tableofcontents
\newpage
\setcounter{section}{4}\setcounter{figure}{2}\setcounter{table}{4}\setcounter{equation}{4}
% ============================== 5 =========================================
\section{Implementación en Python}

\subsection{Estructura del repositorio}

\begin{lstlisting}[language={}, numbers=none]
agente_devops/
  base.py        Plantilla del taller (clase AgenteBasadoEnUtilidad)
  sensores.py    SensoresPR + parsers de git diff, JUnit, coverage y Bandit
  agente.py      AgenteDevOps: modelo del mundo, transicion y actuadores
  kpis.py        Metricas de desempeno
  config.py      Pesos y umbrales calibrados
scripts/
  recolectar_sensores.py   Construye la percepcion en GitHub Actions
  revisar_pr.py            Ejecuta el agente y genera el comentario del PR
simulacion.py    Bateria de 7 casos de prueba + analisis de sensibilidad
tests/           26 pruebas unitarias (pytest), cobertura 97 %
.github/workflows/  ci.yml y agente-pr.yml
\end{lstlisting}

El código conserva la interfaz abstracta de la plantilla y la extiende en dos puntos: la función de
utilidad incorpora los términos $D_q$ y $w_p$, y \texttt{EstadoProyectado} valida que cada variable
esté en su rango. Los parámetros viven en \texttt{config.py}, separados de la lógica, para poder
recalibrar sin tocar el código.

\begin{lstlisting}[caption={Función de utilidad (\texttt{base.py}).}]
def calcular_utilidad(self, estado: EstadoProyectado) -> float:
    p = self.pesos
    utilidad = (
        p["exito"] * estado.probabilidad_exito
        + p["beneficio"] * estado.beneficio_operativo
        - p["riesgo"] * estado.riesgo_impacto
        - p["costo"] * estado.costo_intervencion
        - p["calidad"] * estado.deficit_calidad
        - (p["politica"] if estado.viola_politica else 0.0)
    )
    return round(utilidad, 2)
\end{lstlisting}

\begin{lstlisting}[caption={Modelo de riesgo $\Pdef(s)$ (\texttt{agente.py}).}]
def probabilidad_defecto(self, d):
    r = self.factores_riesgo(d)
    p = config.RIESGO_BASE + sum(config.COEF_RIESGO[k] * v for k, v in r.items())
    return round(_acotar(p), 4)
\end{lstlisting}

\paragraph{Manejo de errores.} Los sensores validan sus rangos (líneas negativas, cobertura fuera de
$[0,1]$, más pruebas fallidas que totales) y lanzan \texttt{ValueError}; el agente rechaza pesos
incompletos; el script de revisión devuelve el código 2 si el JSON de sensores es inválido. Cada
decisión se registra con \texttt{logging} y en \texttt{historial\_decisiones}.

\subsection{Integración con GitHub Actions}

El workflow \texttt{agente-pr.yml} se ejecuta en cada PR y materializa el ciclo PEAS completo:

\begin{enumerate}[leftmargin=*]
\item \textbf{Sensores:} \texttt{git diff --numstat}, \texttt{pytest --junitxml --cov}, \texttt{bandit -f json}.
\item \textbf{Percepción:} \texttt{recolectar\_sensores.py} consolida todo en \texttt{sensores.json}.
\item \textbf{Decisión:} \texttt{revisar\_pr.py} ejecuta el agente y escribe \texttt{decision.json} y \texttt{comentario.md}.
\item \textbf{Actuadores:} publica el comentario con la tabla de utilidades, aplica una etiqueta,
      asigna al arquitecto (variable \texttt{ARQUITECTO\_SENIOR}), fusiona si
      \texttt{AGENTE\_AUTOMERGE=true} y deja el \textit{check} en rojo cuando la decisión es
      \textsc{bloquear} o \textsc{corregir}. Con \textit{branch protection} activa, esto impide el merge.
\end{enumerate}

\noindent Como medida de seguridad, el título del PR (entrada controlada por el usuario) se pasa al
script mediante una variable de entorno y nunca se interpola directamente en el comando, lo que evita
la inyección de comandos en el runner.

% ============================== 6 =========================================
\section{Casos de prueba y métricas}

\subsection{Escenarios}

Se diseñaron siete escenarios. Para cada uno se definió la decisión que tomaría un revisor experto
(oráculo) y si el PR contiene realmente un defecto, con el fin de medir los KPIs.

\begin{table}[H]
\centering
\caption{Matriz de casos de prueba (lecturas de sensores).}
\label{tab:casos}
\small
\begin{tabularx}{\textwidth}{@{}l l >{\raggedright\arraybackslash}X r c r c@{}}
\toprule
\textbf{PR} & \textbf{Tipo} & \textbf{Descripción} & \textbf{Líneas} & \textbf{Pruebas} & \textbf{Cob.} & \textbf{SAST} \\
\midrule
101 & Rutinario & Corrige un typo en \texttt{README.md} & 3 & 42/42 & 85\,\% & --- \\
102 & Rutinario & Exportar reportes a CSV, con pruebas & 150 & 58/58 & 85\,\% & --- \\
103 & Crítico & Nuevo cálculo de descuentos & 200 & 48/60 & 80\,\% & --- \\
104 & Crítico & Conexión a BD con contraseña en el código y SQL concatenado & 80 & 40/40 & 82\,\% & 2 crít. \\
105 & Ambiguo & Migración de sesiones a OAuth2 en \texttt{src/auth} & 1200 & 130/130 & 82\,\% & --- \\
106 & Ambiguo & Plantillas de correo sin pruebas suficientes & 300 & 35/35 & 55\,\% & --- \\
107 & Ambiguo & Carga de archivos con \texttt{subprocess(shell=True)} & 60 & 50/50 & 78\,\% & 1 alta \\
\bottomrule
\end{tabularx}
\end{table}

\subsection{Resultados}

\begin{table}[H]
\centering
\caption{Utilidad calculada para cada acción y decisión del agente.}
\label{tab:resultados}
\footnotesize\setlength{\tabcolsep}{4pt}
\begin{tabular}{@{}l l c r r r r l c@{}}
\toprule
\textbf{PR} & \textbf{Diagnóstico} & $\Pdef$ & \textbf{Merge} & \textbf{Corregir} & \textbf{Arquit.} & \textbf{Bloquear} & \textbf{Decisión} & \textbf{Oráculo} \\
\midrule
101 & TRIVIAL    & 0.050 & \textbf{183.00} & 79.00 & 97.00 & 12.00 & Merge & \checkmark \\
102 & ESTANDAR   & 0.072 & \textbf{175.35} & 78.12 & 96.56 & 12.00 & Merge & \checkmark \\
103 & FALLIDO    & 0.170 & $-857.80$ & \textbf{114.20} & 94.60 & 12.00 & Corregir & \checkmark \\
104 & BLOQUEANTE & 0.462 & $-1157.08$ & $-97.48$ & $-111.24$ & \textbf{81.00} & Bloquear & \checkmark \\
105 & RIESGOSO   & 0.300 & 50.00 & 69.00 & \textbf{92.00} & 12.00 & Arquitecto & \checkmark \\
106 & RIESGOSO   & 0.142 & 87.48 & \textbf{115.32} & 95.16 & 12.00 & Corregir & \checkmark \\
107 & RIESGOSO   & 0.263 & $-39.31$ & \textbf{110.48} & 92.76 & 12.00 & Corregir & \checkmark \\
\bottomrule
\end{tabular}
\end{table}

\begin{itemize}[leftmargin=*]
\item \textbf{Rutinarios (101, 102):} el riesgo es bajo y el merge obtiene casi el doble de utilidad
      que cualquier alternativa; el agente libera a los revisores de este trabajo.
\item \textbf{Críticos (103, 104):} la penalización $w_p$ hace imposible el merge. En el PR-103 el
      fallo es concreto y corregible, así que se solicitan correcciones; en el PR-104 la exposición
      $E(s)$ hace que solo bloquear tenga utilidad positiva.
\item \textbf{Ambiguos (105--107):} el PR-105 pasa todas las pruebas, pero su tamaño y criticidad lo
      envían a un arquitecto. En los PR-106 y 107 el término $D_q$ reduce la utilidad del merge por
      debajo de \textsc{corregir}, porque hay hallazgos concretos que el autor puede resolver.
\end{itemize}

\subsection{Indicadores clave de desempeño (KPIs)}

\begin{table}[H]
\centering
\caption{KPIs obtenidos en la simulación.}
\label{tab:kpis}
\begin{tabularx}{\textwidth}{@{}l r X@{}}
\toprule
\textbf{KPI} & \textbf{Valor} & \textbf{Relación con Performance} \\
\midrule
Precisión frente al oráculo experto & 100\,\% & Calidad de las decisiones \\
Defectos escapados a producción & 0 de 4 & Reducción de fallos tras el despliegue \\
\textit{Recall} de defectos & 100\,\% & Todo PR defectuoso fue detenido \\
Tasa de merge automático & 28.6\,\% & Trabajo humano ahorrado \\
Tasa de escalamiento a senior & 14.3\,\% & Solo el caso de alto impacto \\
Tiempo de ciclo promedio & 211.4 min & Frente a 720 min si un senior revisa todo: $-70.6\,\%$ \\
\bottomrule
\end{tabularx}
\end{table}

\noindent El tiempo de ciclo se simuló con valores de referencia por acción: merge 5 min, bloqueo
30 min, corrección del autor 240 min y revisión senior 720 min.

\subsection{Análisis de sensibilidad del peso de riesgo}

Para validar la calibración se repitió la simulación variando $w_r$:

\begin{table}[H]
\centering
\caption{Sensibilidad de las decisiones frente a $w_r$.}
\label{tab:sensibilidad}
\small\setlength{\tabcolsep}{4pt}
\begin{tabular}{@{}r c c c r p{5.2cm}@{}}
\toprule
$w_r$ & \textbf{Precisión} & \textbf{Merge auto} & \textbf{Escapados} & \textbf{Ciclo (min)} & \textbf{Errores} \\
\midrule
50  & 71.4\,\% & 42.9\,\% & 0 & 139.3 & PR-104 no se bloquea; PR-105 se fusiona \\
100 & 85.7\,\% & 42.9\,\% & 0 & 109.3 & PR-105 se fusiona sin revisión senior \\
200 & 85.7\,\% & 42.9\,\% & 0 & 109.3 & PR-105 se fusiona sin revisión senior \\
\textbf{400} & \textbf{100\,\%} & 28.6\,\% & 0 & 211.4 & --- \\
800 & 100\,\% & 28.6\,\% & 0 & 211.4 & --- \\
\bottomrule
\end{tabular}
\end{table}

Con un peso de riesgo bajo, el agente prioriza la velocidad: fusiona sin revisión una migración de
1200 líneas en el módulo de autenticación y, con $w_r=50$, ni siquiera bloquea un secreto expuesto.
A partir de $w_r=400$ las decisiones son estables, por lo que ese valor es el punto en que el agente
resuelve correctamente el dilema sin sacrificar más velocidad de la necesaria. La política dura $w_p$
garantiza que, incluso con pesos mal calibrados, no se fusione código con pruebas rotas.

\subsection{Pruebas unitarias}

La suite de \texttt{pytest} (26 pruebas, 97\,\% de cobertura) verifica los siete escenarios, las
propiedades de seguridad (``nunca fusionar con una vulnerabilidad crítica'' para distintos tamaños de
PR), la monotonía del riesgo frente al tamaño, la fórmula de utilidad, la validación de entradas y
los \textit{parsers} de los artefactos del pipeline. El workflow \texttt{ci.yml} ejecuta la suite, Bandit
y la simulación en cada \textit{push}.

% ============================== 7 =========================================
\section{Conclusiones y trabajo futuro}
\label{sec:conclusiones}

\begin{itemize}[leftmargin=*]
\item El modelo PEAS permitió delimitar desde el inicio qué percibe el agente (artefactos del
      pipeline) y cómo actúa (API de GitHub), lo que se tradujo directamente en el código.
\item La función de utilidad resuelve el dilema velocidad--riesgo de forma cuantitativa y auditable:
      cada comentario en el PR muestra los puntajes de las cuatro acciones.
\item Las restricciones duras ($w_p$) y la exposición $E(s)$ evitan que el equilibrio de pesos
      permita acciones inaceptables, como fusionar un secreto expuesto.
\item El agente resolvió sin intervención humana 6 de los 7 PRs simulados (85.7\,\%: merge,
      corrección o bloqueo) y solo escaló el caso de alto impacto, reduciendo el tiempo de ciclo un 70.6\,\%.
\end{itemize}

\paragraph{Trabajo futuro.}
(1) Evolucionar hacia una arquitectura híbrida, reemplazando los coeficientes $\alpha_i$ de
$\Pdef$ por un modelo de aprendizaje automático entrenado con el historial de incidentes del
repositorio, y usar como elemento de crítica la tasa de reversiones (\textit{reverts}) tras el
despliegue.
(2) Calcular la cobertura solo sobre las líneas modificadas (\textit{diff coverage}).
(3) Incorporar un sistema multiagente (analista de seguridad, crítico de rendimiento y gerente de
merge) orquestado con LangGraph, como propone el ejercicio 3 del taller.

\begin{thebibliography}{9}
\bibitem{russell} S. Russell y P. Norvig, \textit{Artificial Intelligence: A Modern Approach}, 4.ª ed. Pearson, 2021.
\bibitem{taller} \textit{Taller de Ingeniería de Agentes Inteligentes: de los fundamentos a los sistemas multiagente}. Material de clase, noviembre de 2026.
\bibitem{gha} GitHub, ``GitHub Actions documentation''. \url{https://docs.github.com/actions}
\bibitem{bandit} PyCQA, ``Bandit: security linter for Python''. \url{https://bandit.readthedocs.io}
\end{thebibliography}

\end{document}
